\begin{researchbox}{Abstract}
Synthetic LiDAR can expand training data for autonomous-driving perception,
but realistic-looking scans do not necessarily improve performance on real
measurements. This focused narrative review connects general generation
methods to automotive LiDAR and its evaluation, using selected surveys and
primary studies from the project's source collection. Statistical models,
procedural simulation, learned generators and hybrid pipelines offer different
balances of control, data dependence and coverage. LiDAR-specific synthesis
must additionally account for beam geometry, occlusion, missing returns,
intensity and annotation consistency. LiDARsim and SynLiDAR provide evidence
of useful transfer in particular perception settings; LiDARGen demonstrates
scan generation and completion, whose results answer different questions.
The strongest direct evidence of training utility is task performance on
held-out real data under comparable training conditions. Geometry and
distributional metrics help diagnose discrepancies but cannot replace this
assessment. We recommend connecting sensor diagnostics to class- and
range-specific utility, while recording the real data, labels and computation
required. The selected evidence supports targeted supplementation of real
data, rather than a universal replacement strategy. Incomplete systematic
screening and heterogeneous evaluation protocols limit the generality of
the synthesis.
\end{researchbox}
